\section{Exact local arithmetic}

Assume $(Q,21)=1$ and write $Q=3t+a$, where $a\in\{1,2\}$.  Reducing
\eqref{eq:resonance} modulo three gives $p\equiv a\pmod 3$.  Thus every solution has
the unique parameterization
\begin{equation}
 L_1(k)=3k+a,\qquad L_2(k)=16t+3a-7k.               \label{eq:forms}
\end{equation}
Its determinant is
\begin{equation}
 3(16t+3a)+7a=16Q.                                  \label{eq:determinant}
\end{equation}
The shell inequalities imply
$10Q/7<L_1(k)<8Q/5<L_2(k)<2Q$, so the parameter interval has length
$2Q/35+O(1)$.

For a prime $\ell$, let $\nu_Q(\ell)$ be the number of residues $k\pmod\ell$ for
which $L_1(k)L_2(k)\equiv0\pmod\ell$.

\begin{lemma}[local root law]\label{lem:roots}
For $(Q,21)=1$,
\[
 \nu_Q(\ell)=
 \begin{cases}
 1,&\ell\in\{2,3,7\}\text{ or }\ell\mid Q,\\
 2,&\text{otherwise}.
 \end{cases}
\]
In particular, the pair of forms is admissible.
\end{lemma}

\begin{proof}
Modulo two, both forms in \eqref{eq:forms} equal $k+a$.  Modulo three, $L_1=a$
is nonzero and $L_2$ has one root.  Modulo seven, $L_2$ is constant and nonzero:
three times that constant is $2Q\pmod7$; meanwhile $L_1$ has one root.  Outside
$\{2,3,7\}$, both forms have one root, and the two roots coincide exactly when the
determinant \eqref{eq:determinant} vanishes modulo $\ell$.
\end{proof}

Consequently the singular series has the exact shape
\begin{equation}
 \SSer_{3,7}(Q)=C_{3,7}
 \prod_{\substack{\ell\mid Q\\ \ell\ge5}}
 \frac{\ell-1}{\ell-2},                              \label{eq:series}
\end{equation}
where $C_{3,7}>0$ contains the fixed factors and the convergent generic Euler
product.
